% Source: https://artofproblemsolving.com/community/c3h1191827_formula_madness
% Stable identity: aops-topic-1191827-q01
Let the diagonal go from a bottom corner $O$ to the opposite top corner $T$, and let $Q$ be the bottom corner directly below $T$.

\begin{asy}
usepackage("amsmath");
size(4.5cm);
defaultpen(fontsize(10pt));
real s=0.4; pair P(real x,real y,real z){return (x+s*y,z+s*y);}
draw(P(0,0,0)--P(4,0,0)--P(4,0,2)--P(0,0,2)--cycle); draw(P(4,0,0)--P(4,3,0)--P(4,3,2)--P(4,0,2)); draw(P(0,0,2)--P(0,3,2)--P(4,3,2));
draw(P(0,0,0)--P(0,3,0)--P(4,3,0),dashed); draw(P(0,3,0)--P(0,3,2),dashed); draw(P(0,0,0)--P(4,3,2),red+linewidth(0.9)); draw(P(0,0,0)--P(4,3,0),blue+dashed);
label("$x$",P(2,0,0),S); label("$y$",P(4,1.5,0),SE); label("$z$",P(0,0,1),W); label("$d$",P(2,1.5,1),N);
label("\mbox{}",(min(currentpicture,user=true)+max(currentpicture,user=true))/2,black); // keeps the text after the figure black when LaTeX typesets the labels inline
\end{asy}

On the bottom face, $OQ$ is the diagonal of an $x\times y$ rectangle, so $OQ^2=x^2+y^2$ by the Pythagorean theorem. The edge $QT$ has length $z$ and is perpendicular to the bottom face, hence to $OQ$; so triangle $OQT$ is right-angled at $Q$ and
\[d^2=OT^2=OQ^2+QT^2=x^2+y^2+z^2.\]
